\scanpage{231}
Par Eilenberg-Zilber, l'\deleted{\unsure{co}}homologie de $C_*(A_*\otimes_k B_*)$ (\deleted{deux} \inserted{$A_*,B_*$} $k$-modules $\tfrac12$ simpliciaux \deleted{$C_*$} \inserted{$C_*(A_*\otimes B_*)=A_*\otimes B_*$} complexe de chaînes \deleted{pour} op. bord $\partial=(\deleted{\unsure{\sum}}\partial_n)$ $\partial_n=\sum_0^n\partial_{i,n}^*$, -- c'est donc, par Dold-Puppe, aussi $\pi_*(A_*\otimes B_*)$) se calcule en prenant \deleted{du} \inserted{l'homologie du} produit tensoriel \deleted{des complexes} $C_*(A_*)\otimes C_*(B_*)$, qui est lui-$=$ \inserted{``chaîne-}homotope'', par Dold-Puppe, à \deleted{\unsure{$C_*(A$}} $A_*^!\otimes B_*^!$ (car $A_*^!\xrightarrow[\text{chaîne-homotopisme}]{\simeq} A_*$, $B_*^!\to B_*$); \deleted{lequel} \inserted{or son homologie} peut se calculer par Künneth. Ici $B_*^!=(\Lambda^i\Psi_*)^!\approx k[i]$, donc on \deleted{doit avoir} a
\[
 \underset{\substack{\rotatebox{90}{$=$}\\\pi_{\ell}(A_*\otimes B_*)}}{H_{\ell}(A_*^!\otimes B_*^!)}
 \simeq\underset{\substack{\rotatebox{90}{$\simeq$}\\\pi_{\ell-i}(A_*\otimes B_*)}}{H_{\ell-i}(A_*^!)}
\]

\uline{Cor}. 1.5

Considérons p.ex. le complexe \deleted{des chaînes} \inserted{$\tfrac12$ simplicial} De Rham à puissances divisées
\margin{$\big|$}
\[
 A_*^{\overset{\text{degré total}}{*}\overset{\text{degré ext}}{*}}\simeq\Gamma^*(\Phi_*)\otimes_k\Lambda^*\Psi_*
\]
on trouve \deleted{\unsure{que}} alors
\[
 \pi_\ell(A_*^{*i})\simeq\pi_{\ell-i}(\underbrace{A_*^{*0}}_{\Gamma^*(\Phi_*)})
 \xleftarrow{\sim}\pi_{\ell-i}(\underbrace{A_*^{00}}_{\Gamma^*(0)=k_*})
 =\begin{cases}0&\text{si }\ell\ne i\\ k&\text{si }\ell=i\end{cases}
\]
En \deleted{\unsure{Alors}} particulier \deleted{\unsure{si}}
\[
 \pi_\ell(\underbrace{A_*^{p,i}}_{\Gamma_k^{p-i}\Phi_*\otimes\Lambda^i\Psi_*})
 \simeq\pi_{\ell-i}(\deleted{\unsure{\Gamma}}\underbrace{A_*^{p-\deleted{\unsure{i}}i,0}}_{\Gamma_*^{p-i}\Phi_*})
 \simeq\begin{cases}0&\text{si }\ell\ne i\text{ ou }p\unsure{\ne|=}i\\ k&\text{si }i=\ell=p\end{cases}
\]
donc, pour $p$ fixé, $i$ variable, on trouve que les \inserted{\deleted{\unsure{com|comp.}} composantes} $A_*^{p,0}$, $A_*^{p,1}$ -- $A_*^{p,p-1}$ \deleted{de} (en degré $<p$) de $A_*^{p*}$ \deleted{sont} (qui est une résolution de $k_*$) sont flasques, la seule composante non flasque étant $A_*^{p,p}$ (ayant comme seul gr. d'hom non nul $\pi_p(A_*^{p,p})=k$). On peut donc considérer que

% archive sheet number: 230

